\section{A graph is a pattern of connections}\label{ch07:sec:graph}
Many situations consist of objects, some pairs of which are connected: people and the handshakes between them, towns and the roads that join them, students and the classmates they know. To study such a situation mathematically, we keep only two pieces of information: which objects there are, and which pairs of them are connected. Everything else, such as how long a road is or where a town lies on a map, is set aside. The structure that remains is called a graph.

\par\noindent\begin{minipage}{\linewidth}
\begin{definition}[Graph]\label{ch07:def:graph}
A \emph{graph} consists of a finite set $V$, whose members are called \emph{vertices}, and a set $E$, whose members are called \emph{edges}. Each edge is a set $\{u,v\}$ of two different vertices $u$ and $v$. These two vertices are the \emph{endpoints} of the edge, and we say that the edge \emph{joins} them.
\end{definition}
\end{minipage}\par

The full name of this kind of graph is a \emph{finite simple undirected graph}, and each of the three words records a choice.
\begin{itemize}
\item \emph{Finite}: the set $V$ of vertices is finite, in the sense of Section~\ref{ch03:sec:sets}. The set $E$ of edges is then finite too, because a finite set has only finitely many pairs of members.
\item \emph{Undirected}: an edge has no direction. Since the order in which members of a set are listed does not matter, $\{u,v\}=\{v,u\}$, and the edge joins $u$ to $v$ exactly as much as it joins $v$ to $u$. Section~\ref{ch07:sec:sink} introduces edges that do have a direction.
\item \emph{Simple}: no edge joins a vertex to itself, because the two endpoints of an edge must be different. Such an edge would be called a \emph{loop}. Also, two vertices are joined by at most one edge: $E$ is a set, and listing $\{u,v\}$ twice does not give it a second member, just as $\{1,1,3\}=\{1,3\}$. Some applications need loops or several edges between the same two vertices, but we will not.
\end{itemize}
Throughout this book, the word \emph{graph} on its own means a graph in the sense of Definition~\ref{ch07:def:graph}. Graphs whose edges have a direction are always called \emph{directed} graphs.

We usually write the edge $\{u,v\}$ more briefly as $uv$. Since $\{u,v\}=\{v,u\}$, the names $uv$ and $vu$ denote the same edge. In Chapter~\ref{ch:start}, two letters written side by side meant a product. In this chapter, two vertex names written side by side mean an edge, and the context always makes clear which is intended.

Here is an example that we will use throughout the chapter. Its vertex set is $V=\{a,b,c,d\}$, and its edge set is
\[
E=\{ab,\ bc,\ ca,\ cd\}.
\]
We call it the \emph{example graph} $G$. In a drawing, each vertex is a small circle, and each edge is a line between its two endpoints:
\begin{center}
\begin{tikzpicture}[every node/.style={circle,draw=ink,minimum size=6mm,inner sep=1pt},line width=.7pt]
\node(a)at(0,0){$a$};\node(b)at(2,0){$b$};\node(c)at(1,1.5){$c$};\node(d)at(3,1.5){$d$};
\draw(a)--(b)--(c)--(a);\draw(c)--(d);
\end{tikzpicture}
\end{center}
A drawing is only a picture of a graph. The positions of the circles, the lengths and shapes of the lines, and any points where two lines happen to cross carry no information; a crossing point is not a vertex. If we moved the circle for $d$ to the far left of the page, we would get a different picture of the same graph, because the same pairs of vertices would still be joined. Chapter~\ref{ch:structure} made the matching point about names: two graphs joined by an isomorphism differ only in the names of their vertices, and they have the same pattern of connections.

Two vertices $u$ and $v$ are \emph{neighbors}, or \emph{adjacent}, when $uv$ is an edge. The \emph{degree} of a vertex $v$, written $\deg(v)$, is the number of its neighbors. In the example graph, the neighbors of $c$ are $a$, $b$, and $d$, so $\deg(c)=3$. The four degrees are
\[
\deg(a)=2,\qquad\deg(b)=2,\qquad\deg(c)=3,\qquad\deg(d)=1.
\]
A vertex with no neighbors at all has degree $0$. It still belongs to the graph; it is simply not an endpoint of any edge.

The degree of $v$ can also be read as the number of edges that have $v$ as an endpoint. Each neighbor $w$ of $v$ gives one such edge, namely $vw$. Different neighbors give different edges, and every edge with endpoint $v$ arises in this way from its other endpoint. Both halves of the word ``simple'' are needed here. If loops were allowed, a loop at $v$ would be an edge at $v$ that comes from no neighbor; if two edges could join the same pair of vertices, one neighbor could account for two edges. We will use this second reading of the degree in the next section.

\pauseq{In a graph with $n$ vertices, can a vertex have degree larger than $n-1$?}{No. The neighbors of a vertex $v$ are vertices different from $v$, because the two endpoints of an edge are different. So the neighbors of $v$ lie among the other $n-1$ vertices, and $\deg(v)\le n-1$. This bound is reached exactly when $v$ is joined to every other vertex. In the example graph, which has $4$ vertices, the vertex $c$ has the largest possible degree, $3$.}
